\begin{thebibliography}{99}
\small
\setlength{\itemsep}{1pt}
\setlength{\parskip}{0pt}
\bibitem{ipc} DANE (2026). \textit{IPC. Índices, serie de empalme 2003--2025}. Diciembre de 2025. \href{https://www.dane.gov.co/files/operaciones/IPC/dic2025/anex-IPC-Indices-dic2025.xlsx}{Anexo estadístico}.
\bibitem{productividad} DANE (2026). \textit{Productividad laboral 2005--2025pr}. Cuadros 1 y 2. Publicación del 31 de marzo. \href{https://www.dane.gov.co/files/operaciones/PTF/anex-PTF-ProductividadLaboral-2025.xlsx}{Anexo estadístico}.
\bibitem{salarios} Gobierno de Colombia. Decretos 1785 de 2020, 1724 de 2021, 2613 de 2022, 2292 de 2023 y 1572 de 2024. Montos del salario mínimo de 2021 a 2025. \href{https://cancilleria.gov.co/normograma/compilacion/docs/decreto_2613_2022.htm}{Decreto 2613 de 2022}. \href{https://www.funcionpublica.gov.co/eva/gestornormativo/norma.php?i=257156}{Decreto 1572 de 2024}.
\bibitem{decreto2026} Gobierno de Colombia (2026). \textit{Decreto 159 de 2026}. Artículo 1 y motivación del aumento del salario mínimo. \href{https://cancilleria.gov.co/normograma/compilacion/docs/decreto_0159_2026.htm}{Texto de la norma}.
\bibitem{ipm} Banco de la República, equipo técnico (2026). \textit{Informe de Política Monetaria. Julio de 2026}. Publicado el 4 de agosto. Resumen macroeconómico, p. 7. Proyección de inflación de 6,9\% para diciembre de 2026. \href{https://www.banrep.gov.co/es/publicaciones-investigaciones/informe-politica-monetaria/julio-2026}{Informe}.
\bibitem{bancolombia} Bancolombia (2026). \textit{Bancolombia proyecta que la inflación 2027 cerrará en 5,5\%}. Actualización del 23 de septiembre, sección de inflación y tasas de interés. Pronóstico para diciembre de 2026 de 6,7\%. \href{https://prensa.bancolombia.com/noticias/economia-finanzas/proyecciones-de-inflacion-bancolombia}{Comunicado}.
\bibitem{minutas} Banco de la República (2026). \textit{Minutas de la reunión de septiembre de 2026}. Publicadas el 5 de octubre. Expectativas de los analistas para fin de 2026 de 6,8\%. \href{https://banrep.gov.co/es/noticias/junta-directiva/minutas-banrep-septiembre-2026}{Minutas}.
\bibitem{eof} Fedesarrollo (2026). \textit{Encuesta de Opinión Financiera. Resultados octubre 2026}. Divulgación del 21 de octubre. Se consulta la página oficial de divulgación y su registro de la expectativa de inflación de 6,87\%. \href{https://fedesarrollo.org.co/noticias/encuesta-de-opinion-financiera-resultados-octubre-2026}{Divulgación de resultados}.
\bibitem{ipc2026} DANE (2026). \textit{IPC. Información técnica, septiembre de 2026}. Publicación del 7 de octubre. \href{https://www.dane.gov.co/index.php/estadisticas-por-tema/precios-y-costos/indice-de-precios-al-consumidor-ipc/ipc-informacion-tecnica}{Información técnica}.
\bibitem{ipc2025} DANE (2026). \textit{Comunicado de prensa. IPC, diciembre de 2025}. \href{https://www.dane.gov.co/files/operaciones/IPC/dic2025/cp-IPC-dic2025.pdf}{Comunicado}.
\bibitem{debate} Suárez, M. C. (2026). \textit{Así quedaría el salario mínimo 2027 en Colombia con el aumento del 9\% que pide Fenalco}. Bloomberg Línea, 27 de octubre. Se utiliza como registro del debate público. \href{https://www.bloomberglinea.com/latinoamerica/colombia/asi-quedaria-el-salario-minimo-2027-en-colombia-con-el-aumento-del-9-que-pide-fenalco/}{Artículo}.
\bibitem{oit} Organización Internacional del Trabajo (2025). \textit{Economic factors for wage setting}. DOI 10.54394/GXOE3754. \href{https://www.ilo.org/publications/economic-factors-wage-setting}{Publicación}.
\bibitem{dube} Dube, A. (2019). \textit{Impacts of minimum wages. Review of the international evidence}. Informe independiente para el Gobierno del Reino Unido. \href{https://www.gov.uk/government/publications/impacts-of-minimum-wages-review-of-the-international-evidence}{Informe y síntesis oficial}.
\bibitem{espe} Arango-Thomas, L. E. (coord.) y coautores (2022). \textit{Efectos macroeconómicos del salario mínimo en Colombia}. Ensayos sobre Política Económica, 103. DOI 10.32468/espe103. \href{https://investiga.banrep.gov.co/es/espe103}{Publicación}.
\bibitem{orozcoRevision} Orozco, J. M. (2026a). \textit{Criterios para la fijación del salario mínimo. Revisión comparada de experiencias en el continente americano}. Documento de trabajo, octubre.
\bibitem{orozcoBoletin} Orozco, J. M. (2026b). \textit{Salario mínimo 2027. El debate pendiente}. Borrador del 27 de octubre.
\end{thebibliography}
